% ----------------------
\subsection{Section 86 (6 December 1832)}
% ----------------------
	\label{DC86}
	
	% -----
	\subsubsection{``ON PRIESTHOOD''}
	% -----
	
		In the 1835 and 1844 DC versions of this section, it has a heading above it which reads ``ON PRIESTHOOD.''
	
	% -----
	\subsubsection{Historical Background}
	% -----
		\label{DC86-HB}
		
		% --
		\paragraph{JSP}
		% --
		
			``JS wrote in his journal that on 6 December 1832, he spent part of the day `translating,' or working on his revision of the Bible. It is not known whether he was working that day on the New Testament revision or the Old Testament revision. But on that same day, he `received a Revelation explaining the Parable [of] the wheat and the tears [tares],' found in Matthew 13, suggesting that he may have been working on the New Testament. When JS worked on that parable more than a year earlier while revising the New Testament, he made few significant changes. Between late July 1832 and early February 1833, however, he apparently spent time reviewing his revisions to the New Testament. At some point, JS and Sidney Rigdon changed the text of Matthew 13:30 (which JS had originally left intact) from `I will say to the reapers, gather ye together first the tares' to `gather ye together first the wheat into my barn, and the tares are bound in bundles to be burned.' This inverted order followed the eschatological sequence of events outlined in a November 1831 revelation: the righteous were to `flee unto Zion' and Jerusalem, leaving the wicked nations behind. This 6 December revelation, which has the Lord telling the angels to `first gather out the wheat,' goes in the same direction as that revision, changing the wording slightly regarding the disposition of the tares. Whether the revelation was dictated before or after the revision was made is unclear, as the revision could have been made anywhere within a roughly six-month window of time.

			``It is also possible that the `translating' JS mentioned in his 6 December journal entry referred not to his review of his earlier New Testament revisions but to his work of revising the Old Testament, which he was engaged in at the same time. Between July 1832, when Frederick G. Williams became the principal scribe for JS's revision of the Old Testament, and July 1833, when JS and Williams completed that work, Rigdon filled in as scribe only once---for the revision of Jeremiah 18--24. Since an extant copy of the 6 December revelation attests that Rigdon wrote the revelation as JS dictated it, Rigdon may have helped JS with the Bible revision that day, in which case JS may have been revising those chapters in Jeremiah on 6 December. The chapters include passages on the scattering and gathering of Israel, including verses explaining that the Lord would `gather the remnant of my flock out of all countries whither I have driven them' and `set up shepherds over them which shall feed them.'

			``The revelation on the wheat and the tares emphasizes the gathering of the righteous in the last days. It incorporates elements of the book of Revelation to recast the parable as a history of Christianity from the days of the apostles to the world's end. The description of a second sowing in the last days clearly depicts Mormonism as a restoration of primitive Christianity. Likewise, the end of the revelation expounds on the idea of priesthood, addressing those in `whom the priesthood hath continued through the lineage of your fathers.' The revelation seems to indicate that those ordained to the priesthood are essential to the gathering of Israel, as it counsels them to be a `light unto the Gentiles' and a `savior' to Israel. The priesthood component of the revelation was apparently perceived as its key aspect: upon its publication in the 1835 edition of the Doctrine and Covenants, the revelation bore the heading `On Priesthood.'

			``The original manuscript of this revelation is no longer extant. Frederick G. Williams copied it into Revelation Book 2, probably between late January and late February 1833.''
			
		% --
		\paragraph{HDDC}
		% --
		
			``December 1832 must have been one of the most awesome periods in the life of the Prophet as he viewed future events related to the wars and destruction prior to the second coming. In Section 86 he saw the eventual separation of the wheat and tares as prophesied by the Savior. Later, he had revealed to him the extent of the wars that would cover the earth prior to the second coming of Christ. He was also commanted to write and have published much of what he understood would happen. Accordingly, he wrote Mr. N.E. Seaton, the editor of a newspaper the name of which cannot now be ascertained, and warned...%
				\footnote{%
					% \label{}%
					Here Woodford quotes from a letter JS wrote to Noah Saxton on 4 January 1833; see that letter \hyperref[EMS-1833-JS-4-JAN-1833]{here}.
					}
					
			``Joseph later recorded in his history that he had seen in visions the end of this nation and the breaking up of the government if it continued to disregard the rights of the citizens. Elder Grant also recorded the extent of the visions by Joseph Smith concerning these matters. He wrote: 
			
				\begin{quote}
					The Prophet stood in his own house when he told several of us of the night the visions of heaven were opened to him, in which he saw the American continent drenched in blood, and he saw nation rising up against nation. He also saw the father shed the blood of the son, and the son the blood of the father; the mother put to death the daughter, and the daughter the mother; and natural affection forsook the hearts of the wicked; for he saw that the Spirit of God should be withdrawn from the inhabitants of the earth, in consequence of which there should be blood upon the face of the whole earth, except among the people of the Most High. The Prophet gazed upon the scene his vision presented, until his heart sickened, and he besought the Lord to close it up again (JD 2:147).
				\end{quote}
				
			``Not all of these things were necessarily seen during December 1832, but enough was revealed for Joseph to respond with these two revelations%
				\footnote{%
					% \label{}%
					DC 86 and 87.
					}
			and the letter cited above, plus a portion of a later revelation (See Section 130).
			
			``Then at the end of the month, as though the Lord sought to comfort Joseph's heart, he revealed Section 88, which Joseph identified as `The Olive Leaf,' plucked from the Tree of Paradise, or the Lord's mesage of peace to His saints. Thus hope was given the members of the Church that they could escape much of the future turmoil if they would live righteously'' (1093--1095).
			
	% -----
	\subsubsection{Versions}
	% -----
		\label{DC86-V}
		
		See the \href{https://drive.google.com/file/d/1goAvNz8jS4R8slYfMG_db55Ty2z_vmgK/view?usp=sharing}{sections comparison} document.
	
		\begin{compactdesc}
			\item[JSP] \hyperref[EMS-1832-JS-6-DEC-1832]{6 December 1832}
			\item[BC] ---
			\item[DC 1835] Section 6, 99
			\item[TS] 5:20 (1 November 1844), 688
			\item[DC 1844] Section 6, 128--129
			\item[MS] 14:19 (3 July 1852), 295--296
		\end{compactdesc}

	% -----
	\subsubsection{Section 86:1} % *DC86-1 
	% -----
		\label{DC86-1}

			\footnote{%
				% \label{}%
				Remember that, in the 1835 DC version of this section, it has the heading ``ON PRIESTHOOD'' in large, bold letters. It has a similar heading in the 1844 DC.
				}%
		VERILY, thus saith the Lord unto you my servants, concerning the parable of the wheat and of the tares:%
			\footnote{%
				% \label{}%
				See \hyperref[MATTHEW-13-24]{Matthew 13:24--30}.
				}

		% \begin{framed}
		% \scriptsize
			% \begin{compactdesc}
				% \item[JSP] 
				% % \item[BC] 
				% % \item[]
			% \end{compactdesc}
		% \end{framed}

	% -----
	\subsubsection{Section 86:2} % *DC86-2 
	% -----
		\label{DC86-2}

		Behold, verily I say, the field was the world, and the apostles were the sowers of the seed;

		% \begin{framed}
		% \scriptsize
			% \begin{compactdesc}
				% \item[JSP] 
				% % \item[BC] 
				% % \item[]
			% \end{compactdesc}
		% \end{framed}

	% -----
	\subsubsection{Section 86:3} % *DC86-3 
	% -----
		\label{DC86-3}

		And after they have fallen asleep the great persecutor of the church, the apostate, the whore, even Babylon, that maketh all nations to drink of her cup, in whose hearts the enemy, even Satan, sitteth to reign---behold he soweth the tares; wherefore, the tares choke the wheat and drive the church into the wilderness.

		% \begin{framed}
		% \scriptsize
			% \begin{compactdesc}
				% \item[JSP] 
				% % \item[BC] 
				% % \item[]
			% \end{compactdesc}
		% \end{framed}

	% -----
	\subsubsection{Section 86:4} % *DC86-4 
	% -----
		\label{DC86-4}

		But behold, in the last days, even now while the Lord is beginning to bring forth the word, and the blade is springing up and is yet tender---

		% \begin{framed}
		% \scriptsize
			% \begin{compactdesc}
				% \item[JSP] 
				% % \item[BC] 
				% % \item[]
			% \end{compactdesc}
		% \end{framed}

	% -----
	\subsubsection{Section 86:5} % *DC86-5 
	% -----
		\label{DC86-5}

		Behold, verily I say unto you, the angels are crying unto the Lord day and night, who are ready and waiting to be sent forth to reap down the fields;

		% \begin{framed}
		% \scriptsize
			% \begin{compactdesc}
				% \item[JSP] 
				% % \item[BC] 
				% % \item[]
			% \end{compactdesc}
		% \end{framed}

	% -----
	\subsubsection{Section 86:6} % *DC86-6 
	% -----
		\label{DC86-6}

		But the Lord saith unto them, pluck not up the tares while the blade is yet tender (for verily your faith is weak), lest you destroy the wheat also.

		% \begin{framed}
		% \scriptsize
			% \begin{compactdesc}
				% \item[JSP] 
				% % \item[BC] 
				% % \item[]
			% \end{compactdesc}
		% \end{framed}

	% -----
	\subsubsection{Section 86:7} % *DC86-7 
	% -----
		\label{DC86-7}

		Therefore, let the wheat and the tares grow together until the harvest is fully ripe; then ye shall first gather out the wheat from among the tares, and after the gathering of the wheat, behold and lo, the tares are bound in bundles, and the field remaineth to be burned.

		% \begin{framed}
		% \scriptsize
			% \begin{compactdesc}
				% \item[JSP] 
				% % \item[BC] 
				% % \item[]
			% \end{compactdesc}
		% \end{framed}

	% -----
	\subsubsection{Section 86:8} % *DC86-8 
	% -----
		\label{DC86-8}

		Therefore, thus saith the Lord unto you,%
			\footnote{%
				% \label{}%
				Who is this ``you''? Who is the Lord speaking to? It seems like he must be referring to the apostles, to whom he gave the parable originally.
				}
		with whom the priesthood hath continued through the lineage of your fathers---%
			\footnote{%
				% \label{}%
				I don't think we need to see this ``lineage'' as actual blood relations; given later ideas about \hyperref[ADOPTION]{adoption}, the ``fathers'' spoken of here could be anyone that you have a certain relationship to in the priesthood.
				}

		\begin{framed}
		\scriptsize
			\begin{compactdesc}
				\item[JSP] therefore thus saith the Lord unto you with whom the priesthood hath continued through the lineage of your fathers,
				% \item[BC] 
				% \item[]
			\end{compactdesc}
		\end{framed}

	% -----
	\subsubsection{Section 86:9} % *DC86-9 
	% -----
		\label{DC86-9}

		For ye are lawful heirs, according to the flesh, and have been hid from the world with Christ in God---

		\begin{framed}
		\scriptsize
			\begin{compactdesc}
				\item[JSP] for ye are lawful heirs according to the flesh and have been hid from the world with christ in God
				% \item[BC] 
				% \item[]
			\end{compactdesc}
		\end{framed}

	% -----
	\subsubsection{Section 86:10} % *DC86-10 
	% -----
		\label{DC86-10}

		Therefore your life and the priesthood have remained, and must needs remain through you and your lineage%
			\footnote{%
				% \label{}%
				Here I think of Peter, James, and John giving the priesthood to JS and Oliver Cowdery.
				}
		until the restoration of all things spoken by the mouths of all the holy prophets since the world began.

		\begin{framed}
		\scriptsize
			\begin{compactdesc}
				\item[JSP] therefore your life, and the Priesthood hath remained and must needs remain through you and your lineage untill the restoration of all things spoken by the mouth of all the holy Prophets since the world began,
				% \item[BC] 
				% \item[]
			\end{compactdesc}
		\end{framed}

	% -----
	\subsubsection{Section 86:11} % *DC86-11 
	% -----
		\label{DC86-11}

		Therefore, blessed are ye if ye continue in my goodness, a light unto the Gentiles, and through this priesthood, a savior unto my people Israel.  The Lord hath said it.  Amen.

		\begin{framed}
		\scriptsize
			\begin{compactdesc}
				\item[JSP] therefore blessed are ye if ye continue in my goodness, a light unto the Gentiles and th[r]ough this Priesthood a saviour unto my people Israel the Lord hath said it
				% \item[BC] 
				% \item[]
			\end{compactdesc}
		\end{framed}